% Section 4 --- Attack 2: PEFO (~2.1 pages, weight 30%)
%
% REVIEW STATUS: §4.1--4.3 complete draft. §4.4--4.5 partially drafted;
% §4.4.2/4.4.4 are PoC/vendor-dependent (populated at Week 10--23).
%
% OPEN: §4.4.1 candidate name = [DEPLOYED_APP] (placeholder until Task 2.1)
% OPEN: §4.4.2 testnet transactions (populated after Milestone 2a, Task 5.2)
% OPEN: §4.4.4 case file (swap-in from templates at Task 9.4)

\section{Attack 2: Partial-Execution Free Option (PEFO)}
\label{sec:pefo}

\subsection{Concrete Instance: Cross-Rollup Intent Filler}
\label{sec:pefo-instance}

We instantiate the PEFO attack through a cross-rollup intent filler protocol
deployed on two rollups $R_1, R_2$ sharing an Astria shared sequencer
$\mathcal{S}$. The architecture is representative of the class of applications
that rely on shared-sequencer ordering without providing execution atomicity.

\paragraph{Application model.}
A user posts an intent $\mathcal{I} = (\text{sell } a_1 \text{ on } R_1,
\text{receive } a_2 \text{ on } R_2)$: sell asset $a_1$ on rollup $R_1$ and
receive asset $a_2$ on rollup $R_2$. A filler $\mathcal{F}$ monitors both
rollup mempools, identifies profitable intents, and submits fill transactions:
$\tau_1$ (take the sell on $R_1$) and $\tau_2$ (deliver the buy on $R_2$).
Both $\tau_1, \tau_2$ enter their respective Composers, which forward them
to the shared Astria sequencer mempool. The sequencer includes both in a single
\emph{meta block} --- a joint commitment $\Gamma(\tau_1, \tau_2)$ over both
rollup namespaces.

\paragraph{Assumption violated.}
The filler implicitly assumes that if both $\tau_1$ and $\tau_2$ appear in the
same meta block, both will execute successfully. This assumption is precisely
the A2 claim that Theorem~\ref{thm:a2-necessary} shows cannot be provided by
a lazy sequencer. The meta block commits to \emph{ordering}, not execution
outcomes.

%----------------------------------------------------------------------
\subsection{Structural Conditions for PEFO}
\label{sec:pefo-structural}

The PEFO attack is possible whenever three structural conditions hold
simultaneously.

\begin{enumerate}
  \item[\textbf{C1.}] \textbf{Shared commitment over independent STFs.}
    The sequencer $\mathcal{S}$ issues a commitment $\Gamma(B)$ covering
    bundle $B = (\tau_1, \tau_2)$ where $\tau_i$ executes on rollup $R_i$
    with independent state transition function $\delta_i$
    (Definition~\ref{def:independent}).

  \item[\textbf{C2.}] \textbf{Independent execution legs.}
    The legs $\tau_1, \tau_2$ are not connected by a protocol-level execution
    dependency. In particular, there is no cross-chain message (LayerZero
    \texttt{lzReceive}, Wormhole \texttt{receiveWormholeMessages}, etc.)
    that makes $\tau_2$ contingent on $\tau_1$'s execution outcome.

  \item[\textbf{C3.}] \textbf{Lazy sequencer.}
    $\mathcal{S}$ is lazy (Definition~\ref{def:lazy}): $\Gamma(B)$ is issued
    before $\delta_1(\sigma_1, \tau_1)$ or $\delta_2(\sigma_2, \tau_2)$ are
    evaluated.
\end{enumerate}

\noindent When C1--C3 hold, the adversary can construct a free option:

\begin{definition}[PEFO Payoff]
\label{def:pefo-payoff}
The PEFO payoff for an adversary $\mathcal{A}$ holding bundle $B$ under
commitment $\Gamma(B)$ is:
\[
  \Pi^{\mathrm{PEFO}}(B) \;=\; \max\!\left(
    \mathbb{E}[\delta_1(\sigma_1, \tau_1)]
    - \mathbb{E}[\delta_2(\sigma_2, \tau_2)],\; 0
  \right)
\]
where $\mathbb{E}[\delta_i(\sigma_i, \tau_i)]$ denotes the expected value
of executing leg $i$. If leg 1 becomes profitable after $\Gamma(B)$ is issued
but before leg 2 executes, $\mathcal{A}$ can selectively cause leg 2 to fail
(by front-running the resource $\tau_2$ depends on) while allowing leg 1 to
succeed. The payoff is $\Pi^{\mathrm{PEFO}} > 0$ whenever the expected value
of leg 1 exceeds the expected value of leg 2 after the commitment.
\end{definition}

\noindent Conditions C1--C3 are satisfied by any lazy shared sequencer over
independent rollup STFs. We confirm they are satisfied in our primary
demonstration platform (Astria) by direct documentation evidence:

\begin{itemize}
  \item C1: \emph{``a single meta block consisting of transactions submitted
    to its mempool by one or more rollups''}~\cite{Astria-TxFlow}
  \item C3: \emph{``provides a guarantee on the ordering of transactions in
    a block, but it doesn't execute the state transition function (STF) of
    any given rollup''}~\cite{Astria-TxFlow}
\end{itemize}

\noindent C2 is application-level: the intent filler submits independent
fill transactions to each rollup; no cross-chain message relay is in the
critical path.

%----------------------------------------------------------------------
\subsection{PEFO as Theorem 1 Negation}
\label{sec:pefo-negation}

PEFO directly instantiates the contrapositive of Theorem~\ref{thm:a2-necessary}.
In our construction:

\begin{itemize}
  \item The sequencer provides neither capability (i) nor capability (ii):
    it issues $\Gamma(B)$ without a success predicate for either STF (no (i)),
    and has no authority to revert $R_1$'s state after $\tau_1$ succeeds
    (no (ii)).
  \item By Theorem~\ref{thm:a2-necessary}, A2 is not guaranteed.
  \item The PEFO attack constructs the execution scenario from the proof:
    the adversary submits a conflicting transaction $\tau_2'$ that causes
    $\delta_2(\sigma_2, \tau_2)$ to fail while $\delta_1(\sigma_1, \tau_1)$
    succeeds, yielding $\Pi^{\mathrm{PEFO}} > 0$.
\end{itemize}

\paragraph{Relation to existing cross-rollup protocols.}
Three existing protocols close the PEFO gap via capability (ii):

\begin{itemize}
  \item \textbf{Across Protocol} uses UMA oracle + escrow: origin-chain funds
    are held until UMA confirms the destination fill, preventing partial
    finality of the origin leg~\cite{Across-Protocol}.
  \item \textbf{HTLC} uses a hashlock: $\tau_2$ can only execute if
    $\tau_1$'s preimage is revealed, creating a sequential dependency that
    prevents partial execution.
  \item \textbf{UniswapX} signed orders require the filler to deliver the
    exact output amount or the transaction reverts (within-chain atomicity),
    providing a per-chain capability (ii)~\cite{UniswapX}.
\end{itemize}

\noindent None of the above use a shared sequencer in the sense of C1
(they use independent sequencers per chain); they achieve A2 at the
application level via capability (ii), consistent with Theorem~\ref{thm:a2-necessary}.
The §7 design-space map classifies all three.

%----------------------------------------------------------------------
\subsection{Survey of Production Cross-Rollup Applications}
\label{sec:pefo-survey}

\paragraph{Systematic survey.}
We conducted a systematic survey of $N=30$ cross-rollup applications against
the three PEFO structural conditions (C1--C3, Appendix~\ref{app:survey}).
The mandatory filter requires: (TF-1) the application spans $\geq 2$ rollups
under a shared sequencer; (TF-2) bundle legs execute independently without
messaging relay; and (TF-3) the sequencer is lazy.

\paragraph{Structural finding: TF-2 is the binding constraint.}
No deployed application in our survey satisfies TF-1$\wedge$TF-2$\wedge$TF-3
simultaneously. The binding constraint is TF-2: every production cross-rollup
application that satisfies TF-1 and TF-3 routes its cross-rollup operations
through a messaging relay (Hyperlane, LayerZero, or equivalent), which
implements capability~(ii) at the application layer (\S\ref{sec:designspace}).

This is a structural finding, not a search gap. The production ecosystem has
independently converged on Theorem~\ref{thm:a2-necessary}'s necessary
conditions: every cross-rollup application that could face PEFO vulnerability
has added application-level relay coordination (capability ii) or oracle
escrow (capability ii). The gap remains live as a \emph{forward-looking risk}
for applications that build on lazy-sequencer infrastructure without these
protections, as demonstrated by the §4.1 devnet testbed (\S\ref{sec:pefo-instance}).

\paragraph{Coordinated disclosure.}
We disclosed the structural PEFO finding to the shared sequencer vendors
whose platforms satisfy C1 and C3 (see \S\ref{sec:disclosure}). Vendor
responses are documented in Table~\ref{tab:vendor-responses}.

\paragraph{Design-space implication.}
The survey confirms that the design-space map in \S\ref{sec:designspace} is
empirically accurate: every production cross-rollup mechanism instantiates
capability~(i), capability~(ii), or A4-only. No production application
relies on shared-sequencer ordering without adding one of these protections.

%----------------------------------------------------------------------
\subsection{What This Section Claims and Does Not Claim}
\label{sec:pefo-scope}

\paragraph{Claims.}
\begin{enumerate}
  \item The PEFO structural conditions (C1--C3) are satisfied by any lazy
    shared sequencer over independent rollup STFs, confirmed by primary
    source documentation of deployed platforms (Astria, Espresso).
  \item The PEFO attack is executable: the Astria devnet testbed
    (\S\ref{sec:pefo-instance}) demonstrates end-to-end value extraction
    under the structural conditions, with a measured 2,199\,ms PEFO window.
  \item The $N=30$ production survey confirms that deployed systems
    independently converge on Theorem~\ref{thm:a2-necessary}'s mitigations,
    validating the design-space map (\S\ref{sec:designspace}).
\end{enumerate}

\paragraph{Does not claim.}
\begin{enumerate}
  \item This section does not claim that any production application has been
    actively exploited. We demonstrate structural exploitability on a devnet
    testbed and report on the production survey findings.
  \item This section does not characterize the magnitude of real-world losses.
    The \S\ref{sec:meas} measurement study bounds option value empirically,
    but actual loss depends on attacker capital and market conditions.
  \item The absence of a TF-2-satisfying production application in our survey
    is a finding about the current ecosystem state, not a claim that the
    structural gap is permanently unexploitable. As cross-rollup intent
    protocols mature on lazy-sequencer platforms, the gap will become
    immediately relevant.
\end{enumerate}
